\documentclass[a4paper,12pt]{article}
\usepackage{graphicx, setspace, array, xcolor, mathtools, contour, tikz, amsthm, setspace, amsmath,amsfonts,amssymb, subcaption, fancyhdr, lipsum,multicol, geometry, gensymb}
\usepackage{colortbl}
\geometry{a4paper, margin=0.75in}
\contourlength{0.5pt}
\usetikzlibrary{patterns}
\usepackage[english]{babel}
\renewcommand{\qedsymbol}{$\blacksquare$}
\definecolor{darkgreen}{rgb}{0.0, 0.5, 0.0}
\pagestyle{fancy}
\fancyhf{}
\renewcommand{\headrulewidth}{0.4pt}
\fancyhead[L]{\rightmark}
\fancyhead[R]{\thepage}
\title{Modern Algebra}
\author{Assignment 9\\ \\ Yousef A. Abood\\ \\ ID: 900248250}
\date{November 2025}
\setlength{\parindent}{0pt}
\singlespacing
\parskip=1mm
\setcounter{section}{7}
\setcounter{subsection}{1}
\onehalfspacing
\begin{document}
\maketitle
\noindent\makebox[\linewidth]{\rule{15cm}{0.4pt}}
\section{Direct Products}
\subsection*{Problem 4}
\begin{proof}
    Suppose we have two abelien groups $G,H$. To show that $G \oplus H$, pick elements $g,g' \in G$ and $h,h' \in H$. Observe that:
    \begin{align*}
        (g,h)(g',h')=(gg',hh')=(g'g,h'h)=(g',h')(g,h).
    \end{align*}
    Hence, $G \oplus H$ is abelien.
\end{proof}
\subsection*{Problem 5}
\begin{proof}
    Ovserve the direct product $\mathbb{Z}\oplus \mathbb{Z}.$ Assume it is cyclic and choose a pair generator $(a,b)$. That implies that there exists $k \in \mathbb{Z}$ such that $k(a,b)=(1,0).$ \\Clearly, $k=a=\pm 1 \wedge b=0.$ similarly, we should have an element $m \in \mathbb{Z}$ such that $m(a,b)=(0,1)$, but $m(a,b)=m(\pm 1, 0)=(\pm m, 0)=(0,1)$ leading to $0=1$, which is contradiction. Thus, $\mathbb{Z} \oplus \mathbb{Z}$ is not cylic.
    For the second part, the proof is valid. See that if we have a group $G$ with more than one element. Assume the group $\mathbb{Z} \oplus G$ is generated by an element $(x,y)$ such that $x \in \mathbb{Z}, y \in G$. Then there exists $t \in \mathbb{Z}$ such that $(tx, y^t)=(1, e_G).$ Then, $t=x=\pm 1, y=e_G.$ Consider the pair $(0, g), g \ne e_G$ and try to generate it. We pick $s \in \mathbb{Z}$ such that $s(x,y)=(0, g)$
    Observe that $s(x,y)=(sx, y^s)=(0,g)$. So, $s=0$, but $y^s=y^0=e_G,$ which is a contradiction to our choice of $y=g\ne e_G$.
\end{proof}
\subsection*{Problem 6}
We know that $Z_8 \oplus Z_2$ has a subgroup isomorphic to $Z_8$, so there exists an element of order $8.$ However, the group $Z_4\oplus Z_4$ cannot have an element of order $8$, as the possible orders of an element in $Z_4$ are $1, 2, 4.$ Observe the pair $(a,b) \in Z_4 \oplus Z_4$, its order is $\operatorname{lcm}(|a|, |b|)$. For the $\operatorname{lcm}(|a|,|b|)$ to be $8$ then $|a|,|b|=1,8$ or $8,1$ and we cannot find an element of order $8$ in $Z_4.$
\subsection*{Problem 15}
\begin{proof}
    Consider the function $\phi((a,b))=a+bi$ from $\mathbb{R}\oplus \mathbb{R} \to \mathbb{C}$ where $a,b \in \mathbb{R}, a+bi \in \mathbb{C}.$ We start by showing it is a homomorphism, see that:
    \begin{align*}
        \phi((a,b))+\phi((x,y))=(a+bi)+(x+yi)=(a+x)+(b+y)i=\phi((a+x, b+y)).
    \end{align*}
    Next, assume $\phi((a,b))=\phi((x,y))$, so $a+bi=x+yi$. Clearly, $a=x, b=y, (a,b)=(x,y)$ and the function is injective.
    Surjectivity is trivial, observe that for every values $a,b$ in $a+bi$ there is a pair $(a,b)$ in $\mathbb{R} \oplus \mathbb{R}$ such that $\phi((a,b))=a+bi.$
\end{proof}
\subsection*{Problem 17}
\begin{proof}
    We know that any subgroup from a cyclic group is cyclic, and we also know that if we have a direct product $G\oplus H$, then we have a subgroup isomorphic to $G$ and a subgroup isomorphic to $H$. Since any subgroup must be cyclic and $H,G$ are isomorphic to subgroups of the direct product, then $G,H$ are both cyclic. The general case follows, observe the direct product $G_1\oplus G_2\oplus G_3\oplus \cdots \oplus G_i$. We must find a cyclic subgroup of the direct product isomorphic to $G_j, 1\le j \le i.$    
\end{proof}
\subsection*{Problem 21}
For the L.H.S., the order of the group generated by the pair $(g,h)$ is $\operatorname{lcm}(|g|, |h|).$ The order of the R.H.S. is $|g|*|h|.$ Since $\operatorname{lcm}(|g|,|h|)*\gcd(|g|, |h|)=|g|*|h|$, if $\langle (g,h) \rangle = \langle g \rangle \langle h \rangle$, then $|\langle (g,h) \rangle|=|\langle g \rangle||\langle h \rangle|$. That is possible only when $\operatorname{lcm}(|g|,|h|)=|g||h|$, so, $\gcd(|g|, |h|)$ is $1$. Thus, $\langle (g,h) \rangle=\langle g \rangle \langle h \rangle$ when the orders of $g,h$ are relatively primes.
\subsection*{Problem 23}
We investigate the number of elements of order $3$ and construct subgroups of order $3$. We know that there is one identity element in $\mathbb{Z}_3 \oplus \mathbb{Z}_3.$ Since $3$ is prime, then every other element is of order $3$. Thus, any subgroup of order $3$ will consist of $2$ other elements than the identity. Observe the formula:
\[\text{Number of subgroups}= \frac{\text{Number of elements}-1}{2} = \frac{3^n-1}{2}.\]
See that we divide by $2$ to avoid overlapping, as we choose a unique pair of elements in each group.
\subsection*{Problem 35}
We observe that the elements of order $2$ in $\mathbb{C}^*$, that is, $z^2=1, z\ne 1$, then $z=-1.$ Now, observe the group $\mathbb{R}^* \oplus \mathbb{R}^*$ has $3$ elements of order $2$, which are $(-1,1), (-1,-1), (1, -1).$ Clearly, the two groups have different structure, so they are not isomorphic.
\subsection*{Problem 46}
$\mathbb{Z} \oplus D_4 \oplus A_4.$
\subsection*{Problem 48}
The group $Z_2 \oplus Z_2$ has $4$ elements, that is, three elements other than the identity. Since we cannot have an element of order $4$, as if we have, the group would be cyclic. We already know that $Z_2 \oplus Z_2$ is not cyclic as $\gcd(|Z_2|,|Z_2|)\ne 1.$ 
Then, the rest elements are of order $2$. To construct an isomorphism, we map elements of the same order, that is, we have three mapping options for the first element other than the identity, two options for the second element, and one for the remaining one. Thus, we have $6$ automorphisms and we see that each automorphism is a permutation of three elements. Hence, $\operatorname{Aut}(Z_2\oplus Z_2) \cong S_3.$
\end{document}